\chapter*{Appendix I}

\addcontentsline{toc}{chapter}{Appendix I}

\def\shorttitle{Appendix I}

\section*{Options for entering TIC, alkalinity and pH in PHREEQC}
There are a number of options available in PHREEQC to enter the carbon concentration, alkalinity
and the pH. The best option depends on the problem at hand. The options are:

\begin{enumerate}
\item Enter pH and alkalinity. Total moles of carbon(4) is calculated by the program.
\item Enter pH and total carbon. Alkalinity is calculated by the program.
\item Enter alkalinity and total carbon. pH is calculated by the program.
\item Enter pH, estimate of total carbon, and a partial pressure of \co. Program will adjust total carbon
until desired $P_{\co}$ is attained.
\item Enter total alkalinity or total carbon, estimate of pH, and a partial pressure of \co. If possible, the
program will adjust the pH to attain the desired $P_{\co}$.
\end{enumerate}

The first option is used when both pH and alkalinity were determined in the field. The second option
should be used when alkalinity was not titrated in the field but pH and total carbon are known from
the laboratory analysis. The third option can be used to check the measured pH of a water sample in
which both alkalinity and total carbon were analysed. Options 4 and 5 are useful when the \co\
pressure of the water sample is known.